\chapter{Fundamentação teórica}\label{cap:fundamentacao}

A informatização do uso de uma sala de computadores exige mais do que a implementação de formulários. É necessário compreender as atividades realizadas no ambiente, especificar regras verificáveis, preservar a consistência dos dados e apresentar operações de maneira acessível. Este capítulo reúne os conceitos que fundamentaram o desenvolvimento, aproveitando a revisão bibliográfica do projeto de estágio.

\section{Engenharia de Software}\label{sec:engenharia}

A Engenharia de Software reúne processos, métodos e práticas voltados à especificação, construção, verificação e evolução de sistemas. Conforme \citeonline{sommerville2016}, o desenvolvimento deve considerar tanto as funcionalidades necessárias quanto atributos de qualidade, manutenção e adequação ao contexto em que o produto será utilizado. Em problemas institucionais, a análise do processo existente ajuda a evitar que decisões técnicas sejam tomadas sem relação com as necessidades dos usuários.

No estágio, a elaboração de uma solução para a biblioteca começou pelo levantamento do funcionamento manual da sala e pela identificação das tarefas que precisariam ser representadas digitalmente. A definição do escopo foi importante porque o conjunto inicial de possibilidades ultrapassava o que seria desenvolvido na primeira versão funcional.

\subsection{Engenharia de requisitos}

Os requisitos funcionais descrevem os serviços e comportamentos esperados de um sistema, enquanto os não funcionais expressam propriedades e restrições relevantes para seu funcionamento. A ISO/IEC/IEEE 29148:2018 orienta os processos de Engenharia de Requisitos e a elaboração de especificações que possam ser compreendidas, avaliadas e verificadas \cite{iso29148_2018}.

No caso estudado, consultar horários, reservar um computador, registrar entrada e saída e gerar relatórios são exemplos de comportamentos funcionais. Integridade dos registros, clareza das mensagens e acessibilidade da interface exemplificam condições de qualidade. A especificação dessas necessidades deve considerar situações excepcionais, como duas tentativas concorrentes de reserva ou a indisponibilidade de um computador já associado a um atendimento.

Uma regra de negócio é uma restrição ou condição particular do domínio, associada à maneira como uma atividade deve ocorrer. Por exemplo, a consulta de disponibilidade precisa respeitar os horários de funcionamento, e a reserva deve ser rejeitada quando o intervalo solicitado entra em conflito com outra utilização válida. Regras dessa natureza não devem depender apenas da apresentação visual; elas precisam ser aplicadas pela camada responsável pelo processamento da operação.

\subsection{Modelagem e documentação}

A modelagem permite representar atores, atividades e dados antes da implementação, tornando mais explícitas as decisões sobre o funcionamento do produto. Diagramas de casos de uso apresentam objetivos e interações dos participantes com o sistema. Diagramas de atividades descrevem sequências de ações, alternativas e condições de passagem entre etapas. As descrições textuais complementam essas representações ao registrar pré-condições, validações e resultados esperados \cite{sommerville2016}.

Esses artefatos contribuem para a comunicação entre as pessoas envolvidas e para a rastreabilidade das funcionalidades. Quando uma regra é alterada, é possível revisar suas consequências nos fluxos, na implementação e nos testes correspondentes. A documentação, portanto, não se limita ao planejamento: também deve acompanhar as mudanças realizadas durante o desenvolvimento.

\section{Interação Humano-Computador}\label{sec:ihc}

A Interação Humano-Computador (IHC) considera as pessoas, suas tarefas e as condições em que utilizam sistemas interativos. \citeonline{barbosaSilva2010} discutem métodos de identificação de necessidades, elaboração de soluções e avaliação da interação. A qualidade de uma interface não depende somente de sua aparência, mas também da compreensão das informações e da facilidade de concluir tarefas.

A ISO 9241-210:2019 apresenta orientações para o projeto centrado nas pessoas, incluindo compreensão do contexto de uso, envolvimento dos usuários e avaliação iterativa das soluções \cite{iso9241_210_2019}. Para a Sala de Informática, essa abordagem é pertinente porque existem necessidades distintas: quem deseja utilizar um computador precisa consultar disponibilidade e registrar o atendimento; quem acompanha a operação precisa identificar sessões ativas e ocorrências; e quem supervisiona o ambiente precisa configurar parâmetros e consultar indicadores.

A prototipação oferece uma maneira de explorar essas necessidades antes de concluir a implementação. Representações navegáveis das telas permitem discutir a sequência de ações, a distribuição dos elementos e o tratamento de situações excepcionais. Após a implementação, verificações funcionais da navegação e avaliações de usabilidade constituem atividades complementares. A existência de uma interface responsiva ou de testes automatizados, por si só, não comprova que uma avaliação de usabilidade com participantes tenha sido realizada.

\section{Banco de dados e integridade da informação}\label{sec:banco}

Sistemas de banco de dados possibilitam o armazenamento estruturado de informações e a definição de vínculos entre entidades. A modelagem e as restrições de integridade contribuem para impedir estados inconsistentes e assegurar que consultas retornem informações compatíveis com os registros armazenados \cite{silberschatz2019}.

No controle de uso da sala, uma sessão representa a permanência de determinado usuário e pode conter diferentes alocações de computador ao longo do atendimento. Reservas representam compromissos futuros associados a intervalos de tempo, enquanto ocorrências registram problemas relacionados aos equipamentos ou às sessões. Essa separação permite preservar um histórico de trocas sem criar uma sessão distinta para cada computador utilizado.

As operações também precisam considerar concorrência. Duas solicitações simultâneas não devem resultar em reservas confirmadas para o mesmo equipamento e intervalo. Restrições no banco e operações transacionais ajudam a preservar invariantes mesmo quando diferentes solicitações chegam em momentos próximos. O tratamento de integridade deve abranger não apenas a criação de registros, mas também alterações, cancelamentos e correções auditadas.

A dimensão temporal merece atenção particular. Horários regulares de funcionamento, exceções pontuais e versões de turnos podem mudar com o tempo. Para que os relatórios históricos permaneçam compreensíveis, alterações posteriores não devem substituir indiscriminadamente as condições que vigoravam quando uma sessão foi realizada. O armazenamento de referências históricas auxilia na reprodução das totalizações.

\section{Arquitetura de aplicações web}\label{sec:arquitetura}

Uma aplicação web normalmente envolve a interação entre um navegador, componentes responsáveis pelo processamento das requisições e um mecanismo de persistência. A separação de responsabilidades favorece a organização e a manutenção do software \cite{sommerville2016}. Essa separação pode ser implementada de diferentes formas, inclusive em uma única aplicação implantável que mantenha módulos internos bem definidos.

O monólito modular reúne a execução dos componentes em uma aplicação, mas organiza responsabilidades como cadastro, configuração, operações, ocorrências e relatórios em áreas distintas do código. Trata-se de uma alternativa que pode simplificar a implantação de um sistema de escopo delimitado, preservando pontos de evolução. A adequação dessa escolha depende do contexto técnico e operacional, não constituindo uma regra geral para todos os sistemas.

Na arquitetura adotada para a aplicação da biblioteca, o backend centraliza as regras de disponibilidade, reservas e sessões. As páginas web utilizam renderização no servidor e atualizações parciais de conteúdo, enquanto uma API permite expor operações e consultas de forma estruturada. A mesma camada de serviços é reutilizada pelas interfaces, reduzindo o risco de regras diferentes para uma operação semelhante. O detalhamento das tecnologias e módulos empregados é apresentado na Seção~\ref{sec:atividades}.

\section{Acessibilidade, verificação e validação}\label{sec:qualidade}

A acessibilidade digital procura oferecer condições de acesso e interação a pessoas com diferentes necessidades. O Modelo de Acessibilidade em Governo Eletrônico (eMAG), versão 3.1, reúne recomendações pertinentes aos serviços digitais públicos \cite{emag2014}. As Web Content Accessibility Guidelines (WCAG) 2.2 apresentam critérios verificáveis relacionados, entre outros aspectos, à navegação por teclado, à identificação de controles, ao foco visível e ao contraste \cite{wcag22_2024}.

Essas referências orientam a revisão de formulários, mensagens e componentes de navegação. A utilização de recomendações de acessibilidade durante o desenvolvimento deve ser diferenciada da comprovação de conformidade integral, que exigiria uma avaliação sistemática dos critérios pertinentes e o registro dos resultados obtidos.

A verificação busca examinar se o sistema atende às especificações e regras definidas, enquanto a validação procura avaliar sua adequação às necessidades do ambiente de uso \cite{sommerville2016}. Testes de regras, APIs, banco de dados e interfaces podem identificar defeitos e regressões. Já a avaliação em contexto real requer a participação dos responsáveis pelo atendimento e a observação das atividades que o sistema deverá apoiar.

No estágio, o desenvolvimento documentado incluiu verificações automatizadas de funcionalidades e fluxos. Tais evidências sustentam afirmações sobre a implementação testada, mas não autorizam concluir, sem dados adicionais, que houve redução mensurada de tempo de atendimento, ganho de produtividade ou implantação permanente na biblioteca.
